\documentclass[UTF8,a4paper,fontset=windows]{ctexart}
\usepackage[left=9mm,right=9mm,top=8mm,bottom=8mm]{geometry}
\usepackage{amsmath,amssymb,array,booktabs,tabularx,xcolor}
\usepackage[hidelinks]{hyperref}
\xeCJKDeclareCharClass{CJK}{"2460->"2473}
\pagestyle{empty}
\setlength{\parindent}{0pt}
\setlength{\parskip}{2pt}
\setlength{\tabcolsep}{4pt}
\renewcommand{\arraystretch}{1.28}
\newcommand{\Ga}{\operatorname{Ga}}
\newcommand{\IG}{\operatorname{IGa}}
\newcommand{\Be}{\operatorname{Be}}
\newcommand{\Bin}{\operatorname{Bin}}
\newcommand{\Poi}{\operatorname{Pois}}
\newcommand{\Exp}{\operatorname{Exp}}
\newcommand{\E}{\operatorname{E}}
\newcommand{\Var}{\operatorname{Var}}
\newcommand{\MSE}{\operatorname{MSE}}
\newcommand{\BF}{B^*}
\newcommand{\ind}{\mathbf1}
\newcommand{\head}[2]{\noindent{\fontsize{14}{17}\selectfont\bfseries 贝叶斯统计\quad A4 手抄参考}\hfill{\small #1}\par
  \vspace{1pt}\hrule height .7pt\vspace{3pt}{\small #2}\par\vspace{2pt}}
\newcommand{\block}[1]{\par\vspace{3pt}{\bfseries\fontsize{10}{12}\selectfont #1}\par\vspace{1pt}\hrule height .3pt\vspace{3pt}}
\newcommand{\sub}[1]{\par\vspace{2pt}\textbf{#1}\enspace}
\newcommand{\tight}{\setlength{\abovedisplayskip}{2.5pt}\setlength{\belowdisplayskip}{2.5pt}\setlength{\abovedisplayshortskip}{1pt}\setlength{\belowdisplayshortskip}{2pt}}
\newsavebox{\leftbody}
\newsavebox{\rightbody}
\begin{document}
\fontsize{9}{11.1}\selectfont\tight
\head{1 / 分布与共轭先验}{第一章：先验分布与后验分布\hfill 约定：$\Ga(\alpha,\lambda)$ 的 $\lambda$ 为\textbf{率参数}；各小节沿用常见记号。}

\sbox{\leftbody}{\begin{minipage}[t]{.487\textwidth}
\block{1\quad Gamma 函数、Beta 函数与积分}
\[
 \Gamma(\alpha)=\int_0^\infty x^{\alpha-1}e^{-x}\,dx,\quad\alpha>0.
\]
\[
 \Gamma(\alpha+1)=\alpha\Gamma(\alpha),\quad\Gamma(n+1)=n!,\quad\Gamma(\tfrac12)=\sqrt\pi.
\]
\[
 B(\alpha,\beta)=\int_0^1x^{\alpha-1}(1-x)^{\beta-1}\,dx
 =\frac{\Gamma(\alpha)\Gamma(\beta)}{\Gamma(\alpha+\beta)}.
\]
\[
 \int_0^\infty x^{\alpha-1}e^{-\lambda x}\,dx=\frac{\Gamma(\alpha)}{\lambda^\alpha},\qquad
 \int_0^\infty x^{-\alpha-1}e^{-\lambda/x}\,dx=\frac{\Gamma(\alpha)}{\lambda^\alpha}.
\]
$\alpha,\beta,\lambda>0$；$B(1,1)=1$；
$B(\alpha+1,\beta)=\dfrac{\alpha}{\alpha+\beta}B(\alpha,\beta)$。

\block{2\quad 常见分布与常用结论}
\sub{正态 $N(\mu,\sigma^2)$}$x\in\mathbb R$。
\[
 f(x)=\frac{1}{\sqrt{2\pi\sigma^2}}\exp\!\left[-\frac{(x-\mu)^2}{2\sigma^2}\right].
\]
$\E X=\mu$，$\Var X=\sigma^2$，众数、中位数均为 $\mu$；
$\dfrac{X-\mu}{\sigma}\sim N(0,1)$。独立正态之和仍正态，均值、方差分别相加。

\sub{二项 $\Bin(n,\theta)$}$x=0,\ldots,n$。
\[
 p(x\mid\theta)=\binom nx\theta^x(1-\theta)^{n-x},\quad
 \E X=n\theta,\quad\Var X=n\theta(1-\theta).
\]
众数 $\lfloor(n+1)\theta\rfloor$；若 $(n+1)\theta$ 为整数且 $0<\theta<1$，有该数及前一个整数两个众数。$\theta=0,1$ 时分别为 $0,n$。

\sub{均匀 $U(a,b)$}$a<x<b$。
\[
 f(x)=\frac1{b-a},\qquad\E X=\frac{a+b}{2},\quad\Var X=\frac{(b-a)^2}{12}.
\]
区间内密度等高，中位数为 $(a+b)/2$。

\sub{Gamma $\Ga(\alpha,\lambda)$}$x>0$，$\alpha,\lambda>0$；$\lambda$ 为率。
\[
 f(x)=\frac{\lambda^\alpha}{\Gamma(\alpha)}x^{\alpha-1}e^{-\lambda x},\quad
 \E X=\frac{\alpha}{\lambda},\quad\Var X=\frac{\alpha}{\lambda^2}.
\]
众数 $(\alpha-1)/\lambda$（$\alpha>1$）；$\alpha\le1$ 时密度递减。
$cX\sim\Ga(\alpha,\lambda/c)$（$c>0$）。
独立同率 $\Ga(\alpha_i,\lambda)$ 之和为 $\Ga(\sum_i\alpha_i,\lambda)$。

\sub{Beta $\Be(\alpha,\beta)$}$0<x<1$，$\alpha,\beta>0$。
\[
 f(x)=\frac{x^{\alpha-1}(1-x)^{\beta-1}}{B(\alpha,\beta)},\qquad\E X=\frac \alpha{\alpha+\beta}.
\]
\[
 \Var X=\frac{\alpha\beta}{(\alpha+\beta)^2(\alpha+\beta+1)},
\]
\[
 \operatorname{Mode}(X)=\frac{\alpha-1}{\alpha+\beta-2}\quad(\alpha,\beta>1).
\]
$\Be(1,1)=U(0,1)$；内点众数条件不满足时检查 $0,1$ 边界。
$\E(X^r)=B(\alpha+r,\beta)/B(\alpha,\beta)$（$r>-\alpha$）。

\sub{倒 Gamma $\IG(\alpha,\lambda)$}$x>0$，$\alpha,\lambda>0$。
\[
 f(x)=\frac{\lambda^\alpha}{\Gamma(\alpha)}x^{-\alpha-1}e^{-\lambda/x},\qquad
 \operatorname{Mode}(X)=\frac \lambda{\alpha+1}.
\]
\[
 \E X=\frac \lambda{\alpha-1}\ (\alpha>1),\quad
 \Var X=\frac{\lambda^2}{(\alpha-1)^2(\alpha-2)}\ (\alpha>2).
\]
$X^{-1}\sim\Ga(\alpha,\lambda)$；$\E(X^r)=\lambda^r\Gamma(\alpha-r)/\Gamma(\alpha)$（$r<\alpha$）。

\sub{泊松 $\Poi(\lambda)$}$x=0,1,\ldots$，$\lambda>0$。
\[
 P(X=x)=\frac{\lambda^xe^{-\lambda}}{x!},\qquad\E X=\Var X=\lambda.
\]
众数 $\lfloor\lambda\rfloor$；$\lambda$ 为整数时为 $\lambda-1,\lambda$。

\sub{指数 $\Exp(\lambda)=\Ga(1,\lambda)$}$x>0$，$\lambda>0$。
\[
 f(x)=\lambda e^{-\lambda x},\quad\E X=1/\lambda,\quad\Var X=1/\lambda^2.
\]
密度随 $x$ 递减；$P(X>x)=e^{-\lambda x}$（$x\ge0$）。
\end{minipage}}
\sbox{\rightbody}{\begin{minipage}[t]{.487\textwidth}
\block{3\quad 三种信息、似然与后验}
\textbf{总体信息}：总体分布或所属分布族的信息。\\
\textbf{样本信息}：从总体抽取的样本所提供的信息。\\
\textbf{先验信息}：观测前的经验、历史资料等；用 $\pi(\theta)$ 表示。
\[
 L(\theta\mid x)=p(x\mid\theta)=\prod_{i=1}^n p(x_i\mid\theta)\quad(\text{独立样本}).
\]
\[
 \pi(\theta\mid x)=\frac{L(\theta\mid x)\pi(\theta)}{\int_\Theta L(\theta\mid x)\pi(\theta)\,d\theta}
 \ \propto\ L(\theta\mid x)\pi(\theta).
\]
\textbf{顺序}：似然 $\to$ 先验密度 $\to$ 相乘整理 $\to$ 识别、归一化。参数离散时，积分换成求和。

\textbf{贝叶斯假设}（成功概率）：$\theta\sim U(0,1)=\Be(1,1)$。

\block{4\quad 常用共轭先验：教材表 1.3.1}
共轭：先验与后验属于同一分布族。以下 $X_i$ 均为独立同分布样本；参数记号在各模型内分别使用。

\sub{① 二项成功概率}$X\mid\theta\sim\Bin(n,\theta)$，观测 $x$。
\[
 \theta\sim\Be(\alpha,\beta)\quad\Longrightarrow\quad
 \theta\mid x\sim\Be(\alpha+x,\beta+n-x).
\]
若各次为 $\Bin(m_i,\theta)$，则用 $n=\sum_i m_i$，$x=\sum_i x_i$。

\sub{② 泊松均值}$X_i\mid\theta\sim\Poi(\theta)$。
\[
 \theta\sim\Ga(\alpha,\lambda)\quad\Longrightarrow\quad
 \theta\mid x\sim\Ga(\alpha+\textstyle\sum_i x_i,\lambda+n).
\]

\sub{③ 指数率（均值的倒数）}$X_i\mid\theta\sim\Exp(\theta)$。
\[
 \theta\sim\Ga(\alpha,\lambda)\quad\Longrightarrow\quad
 \theta\mid x\sim\Ga(\alpha+n,\lambda+\textstyle\sum_i x_i).
\]

\sub{④ 正态均值（方差已知）}$X_i\mid\theta\sim N(\theta,\sigma^2)$。
\[
 \theta\sim N(\mu,\tau^2)\quad\Longrightarrow\quad
 \theta\mid x\sim N(\mu_n,\tau_n^2),
\]
\[
 \mu_n=\frac{\mu/\tau^2+n\bar x/\sigma^2}{1/\tau^2+n/\sigma^2},\qquad
 \tau_n^2=\frac1{1/\tau^2+n/\sigma^2}.
\]

\sub{⑤ 正态方差（均值已知）}$X_i\mid\sigma^2\sim N(\mu,\sigma^2)$。
\[
 \sigma^2\sim\IG(\alpha,\lambda)\quad\Longrightarrow\quad
 \sigma^2\mid x\sim\IG(\alpha+n/2,\lambda+T/2),
\]
其中 $T=\sum_i(x_i-\mu)^2$；$\E(\sigma^2\mid x)=\dfrac{\lambda+T/2}{\alpha+n/2-1}$（$\alpha+n/2>1$）。

\block{5\quad 后验与先验、样本之间的关系}
\textbf{正态均值：精度相加，均值按精度加权}
\[
 \boxed{\ \frac1{\tau_n^2}=\frac1{\tau^2}+\frac n{\sigma^2}\ },\qquad
 \text{精度}=\frac1{\text{方差}}.
\]
\[
 \mu_n=\gamma\mu+(1-\gamma)\bar x,\qquad
 \gamma=\frac{1/\tau^2}{1/\tau^2+n/\sigma^2}.
\]
注意样本\textbf{均值}的方差为 $\sigma^2/n$，其精度为 $n/\sigma^2$。

\textbf{二项成功概率：先验均值与样本频率加权}
\[
 \E(\theta\mid x)=\frac{\alpha+x}{\alpha+\beta+n}
 =\frac{\alpha+\beta}{\alpha+\beta+n}\frac \alpha{\alpha+\beta}+\frac n{\alpha+\beta+n}\frac xn.
\]

\textbf{正态方差：已知均值下的加权式}（$\alpha>1$）
\[
 \E(\sigma^2\mid x)=\frac{\lambda+T/2}{\alpha+n/2-1}
 =\gamma\frac \lambda{\alpha-1}+(1-\gamma)\frac Tn,
\]
\[
 \gamma=\frac{\alpha-1}{\alpha+n/2-1}.
\]
$T/n$ 以\textbf{已知 $\mu$} 为中心；不是以 $\bar x$ 为中心的 $s^2$。

这些均值是\textbf{加权平均}，权重之和为 $1$；$n$ 增大时，样本信息通常占更大权重。
\end{minipage}}
\typeout{FRONT-LEFT-HEIGHT: \the\dimexpr\ht\leftbody+\dp\leftbody\relax}
\typeout{FRONT-RIGHT-HEIGHT: \the\dimexpr\ht\rightbody+\dp\rightbody\relax}
\noindent\usebox{\leftbody}\hfill\usebox{\rightbody}

\newpage
\head{2 / 贝叶斯推断}{第二章：贝叶斯推断\hfill 统一分位数约定：$P(Y\le q_p)=p$；$z_p=\Phi^{-1}(p)$，$\chi^2_p(\nu)$ 为\textbf{左侧概率 $p$} 分位数。}
\sbox{\leftbody}{\begin{minipage}[t]{.487\textwidth}
\block{6\quad 后验估计、损失与 MSE}
\[
 \hat\theta_E=\E(\theta\mid x),\qquad
 \hat\theta_{\rm Mo}\in\arg\max_\theta\pi(\theta\mid x).
\]
连续后验的中位数 $\hat\theta_{\rm Me}$ 满足
\[
 \int_{-\infty}^{\hat\theta_{\rm Me}}\pi(\theta\mid x)\,d\theta=\tfrac12.
\]
\textbf{平方损失} $L(\theta,\hat\theta)=(\theta-\hat\theta)^2$：后验均值最优。
\[
 \MSE(\hat\theta\mid x)=\E[(\theta-\hat\theta)^2\mid x]
 =\Var(\theta\mid x)+[\hat\theta-\E(\theta\mid x)]^2.
\]
\[
 \Var(\theta\mid x)=\E(\theta^2\mid x)-[\E(\theta\mid x)]^2.
\]
最小后验 MSE 为 $\Var(\theta\mid x)$；后验标准差为其平方根。

\textbf{绝对值损失} $L(\theta,\hat\theta)=|\theta-\hat\theta|$：后验中位数最优。

离散后验的中位数满足 $P(\theta\le \hat\theta\mid x)\ge1/2$ 且 $P(\theta\ge \hat\theta\mid x)\ge1/2$。后验为正态时，均值、中位数、众数相同。

\textbf{频率 MSE}（与上面的后验 MSE 区分）：
\[
 \MSE_\theta(\hat\theta)=\E_\theta[(\hat\theta(X)-\theta)^2]
 =\Var_\theta(\hat\theta)+[\E_\theta\hat\theta-\theta]^2.
\]
前者固定 $x$ 对 $\theta\mid x$ 取期望；后者固定 $\theta$ 对样本取期望。

\block{7\quad 常用充分统计量（独立同分布样本）}
\textbf{定义}：给定 $T(X)$ 后，样本的条件分布与 $\theta$ 无关。
\[
 p(x\mid\theta)=g(T(x),\theta)h(x)\quad\Longleftrightarrow\quad T\text{ 充分}.
\]
贝叶斯形式：\(\pi(\theta\mid x)=\pi(\theta\mid T(x))\)。\par
{\renewcommand{\arraystretch}{1.2}
\begin{tabularx}{\linewidth}{@{}p{49mm} X@{}}
\toprule 模型与未知参数 & 充分统计量 \\
\midrule
Bernoulli / 二项，$\theta$ & $\sum_i X_i$ \\
泊松均值 $\theta$ & $\sum_i X_i$ \\
指数率 / Gamma 率（形状已知） & $\sum_i X_i$ \\
$N(\theta,\sigma^2)$，$\sigma^2$ 已知 & $\bar X$ \\
$N(\mu,\sigma^2)$，$\mu$ 已知 & $\sum_i(X_i-\mu)^2$ \\
$N(\mu,\sigma^2)$，$\mu,\sigma^2$ 均未知 & $(\bar X,\sum_i(X_i-\bar X)^2)$ \\
$U(0,\theta)$ & $X_{(n)}=\max_iX_i$ \\
$U(\theta,\theta+c)$，$c$ 已知 & $(X_{(1)},X_{(n)})$ \\
\bottomrule
\end{tabularx}}
\[
 \sum_i(x_i-\theta)^2=\sum_i(x_i-\bar x)^2+n(\bar x-\theta)^2.
\]
$U(0,\theta)$ 中，似然为 $\theta^{-n}\ind_{\{\theta\ge x_{(n)}\}}$（样本非负）；\textbf{支持集不能略去}。
\[
 P(X_{(n)}\le x\mid\theta)=(x/\theta)^n,\quad0<x<\theta.
\]

\block{8\quad 倒 Gamma $\leftrightarrow$ 卡方：可信上限}
若 $\theta\mid x\sim\IG(\alpha,\beta)$，则
\[
 \theta^{-1}\mid x\sim\Ga(\alpha,\beta),\qquad
 \frac{2\beta}{\theta}\,\bigg|\,x\sim\Ga(\alpha,\tfrac12)=\chi^2(2\alpha).
\]
取倒数后不等号反向，因此 $1-\gamma$ 可信上限为
\[
 \boxed{\ \theta_U=\frac{2\beta}{\chi^2_{\gamma}(2\alpha)}\ },\qquad
 P(\theta\le\theta_U\mid x)=1-\gamma.
\]
核查：$\theta\le\theta_U\iff 2\beta/\theta\ge\chi^2_\gamma(2\alpha)$，右侧概率为 $1-\gamma$。
\[
 \theta_L=\frac{2\beta}{\chi^2_{1-\gamma}(2\alpha)},\quad P(\theta\ge\theta_L\mid x)=1-\gamma.
\]
等尾 $1-\gamma$ 可信区间：
\[
 \left[\frac{2\beta}{\chi^2_{1-\gamma/2}(2\alpha)},\ \frac{2\beta}{\chi^2_{\gamma/2}(2\alpha)}\right].
\]
正态后验 $\theta\mid x\sim N(\mu_n,\tau_n^2)$：
\[
 \theta_U=\mu_n+z_{1-\gamma}\tau_n;\qquad\text{等尾区间 }\mu_n\pm z_{1-\gamma/2}\tau_n.
\]
\end{minipage}}
\sbox{\rightbody}{\begin{minipage}[t]{.487\textwidth}
\block{9\quad 假设检验：统一公式与三种情形}
$H_i:\theta\in\Theta_i$，$\Theta_0\cap\Theta_1=\varnothing$，$\Theta_0\cup\Theta_1=\Theta$。

记\textbf{先验概率} $\pi_i=P(H_i)>0$，\textbf{后验概率} $\alpha_i=P(H_i\mid x)$，以及 $m_i(x)=p(x\mid H_i)$。
\[
 \alpha_i=\frac{\pi_i m_i(x)}{\pi_0m_0(x)+\pi_1m_1(x)},\qquad i=0,1.
\]
\[
 \underbrace{\frac{\alpha_0}{\alpha_1}}_{\text{后验机会比}}
 =\underbrace{\frac{\pi_0}{\pi_1}}_{\text{先验机会比}}
 \underbrace{\frac{m_0(x)}{m_1(x)}}_{\BF(x)\text{：支持 }H_0\text{ 的 Bayes 因子}}.
\]
\[
 \BF(x)=\frac{\alpha_0/\alpha_1}{\pi_0/\pi_1}=\frac{m_0(x)}{m_1(x)}.
\]
$\BF>1$：数据支持 $H_0$；$\BF<1$：数据支持 $H_1$。最终选择还取决于先验概率和误判损失。

\sub{① 简单对简单}$H_0:\theta=\theta_0$，$H_1:\theta=\theta_1$。
\[
 m_0(x)=p(x\mid\theta_0),\quad m_1(x)=p(x\mid\theta_1),\quad
 \BF(x)=\frac{p(x\mid\theta_0)}{p(x\mid\theta_1)}.
\]

\sub{② 复合对复合}在 $H_i$ 内取归一化密度 $g_i$：
\[
 g_i(\theta)=\frac{\pi(\theta)\ind_{\Theta_i}(\theta)}{\pi_i},\quad
 \pi_i=\int_{\Theta_i}\pi(\theta)\,d\theta.
\]
\[
 m_i(x)=\int_{\Theta_i}p(x\mid\theta)g_i(\theta)\,d\theta,\qquad
 \BF(x)=\frac{\int_{\Theta_0}p(x\mid\theta)g_0(\theta)\,d\theta}
 {\int_{\Theta_1}p(x\mid\theta)g_1(\theta)\,d\theta}.
\]
也可先求后验，再积分：$\alpha_i=\int_{\Theta_i}\pi(\theta\mid x)\,d\theta$。

\sub{③ 简单对复合}$H_0:\theta=\theta_0$，$H_1:\theta\ne\theta_0$。

必须给点假设\textbf{正的先验质量}：$P(\theta=\theta_0)=\pi_0$；在 $H_1$ 下用密度 $g_1$。纯连续先验给单点的概率为 $0$。
\[
 m_0(x)=p(x\mid\theta_0),\qquad
 m_1(x)=\int_{\Theta_1}p(x\mid\theta)g_1(\theta)\,d\theta.
\]
\[
 \BF(x)=\frac{p(x\mid\theta_0)}{m_1(x)},\qquad
 \alpha_0=\left[1+\frac{1-\pi_0}{\pi_0}\frac1{\BF(x)}\right]^{-1}.
\]
$\pi_0=\pi_1=1/2$ 时，$\alpha_0=\BF/(1+\BF)$，$\alpha_1=1/(1+\BF)$。

\sub{如何作决策}正确决策损失为 $0$；记 $c_{ij}$ 为\textbf{选择 $H_i$、实际 $H_j$} 的损失，则
\[
 R(H_0\mid x)=c_{01}\alpha_1,\quad R(H_1\mid x)=c_{10}\alpha_0.
\]
选择 $H_0\iff c_{01}\alpha_1<c_{10}\alpha_0$；相等时任取。等误判损失时，选后验概率较大的假设；即选 $H_0\iff\alpha_0>1/2$。

\block{10\quad 预测：先求后验，再对参数积分}
设未来量 $Z$ 的条件密度 / 分布列为 $g(z\mid\theta)$，给定 $\theta$ 后 $Z$ 与已有样本 $X$ 条件独立，则
\[
 \boxed{\ m(z\mid x)=\int_\Theta g(z\mid\theta)\pi(\theta\mid x)\,d\theta\ }.
\]
\textbf{步骤}：①求 $\pi(\theta\mid x)$；②写 $g(z\mid\theta)$；③相乘，对 $\theta$ 积分（离散参数求和）；④保留 $z$ 的支持集。
\[
 P(Z\in C\mid x)=\int_\Theta P(Z\in C\mid\theta)\pi(\theta\mid x)\,d\theta.
\]
积分后得到的是\textbf{未来观测的分布}，不是参数后验；一般不等于把后验均值代进 $g(z\mid\theta)$。
\end{minipage}}
\typeout{LEFT-COLUMN-HEIGHT: \the\dimexpr\ht\leftbody+\dp\leftbody\relax}
\typeout{RIGHT-COLUMN-HEIGHT: \the\dimexpr\ht\rightbody+\dp\rightbody\relax}
\noindent\usebox{\leftbody}\hfill\usebox{\rightbody}

\newpage
\head{3 / 教材例 1.5.1}{正态均值与方差均未知：联合共轭模型\hfill 沿用教材和板书的 $\mu,\sigma^2,\mu_0,\kappa_0,\nu_0,\sigma_0^2$。}
\noindent\begin{minipage}[t]{.487\textwidth}
\block{11\quad 模型、似然与联合先验}
设 $X_1,\ldots,X_n\mid\mu,\sigma^2\overset{\rm iid}{\sim}N(\mu,\sigma^2)$，
$\mu,\sigma^2$ 均未知；$\kappa_0,\nu_0,\sigma_0^2>0$。
\[
 \bar x=\frac1n\sum_i x_i,\qquad(n-1)s^2=\sum_i(x_i-\bar x)^2.
\]
\textbf{Step 1：似然函数}
\[
 \begin{aligned}
 L(\mu,\sigma^2\mid x)
 &=(2\pi\sigma^2)^{-n/2}\exp\!\left[-\frac{\sum_i(x_i-\mu)^2}{2\sigma^2}\right]\\
 &\propto(\sigma^2)^{-n/2}\exp\!\left[-\frac{(n-1)s^2+n(\bar x-\mu)^2}{2\sigma^2}\right].
 \end{aligned}
\]
\textbf{Step 2：先验密度（正态—倒 Gamma）}
\[
 \mu\mid\sigma^2\sim N\!\left(\mu_0,\frac{\sigma^2}{\kappa_0}\right),\qquad
 \sigma^2\sim\IG\!\left(\frac{\nu_0}{2},\frac{\nu_0\sigma_0^2}{2}\right).
\]
\[
 \pi(\mu,\sigma^2)=\pi(\mu\mid\sigma^2)\pi(\sigma^2).
\]
\[
 \pi(\mu\mid\sigma^2)=\sqrt{\frac{\kappa_0}{2\pi\sigma^2}}
 \exp\!\left[-\frac{\kappa_0(\mu-\mu_0)^2}{2\sigma^2}\right].
\]
\[
 \pi(\sigma^2)=\frac{(\nu_0\sigma_0^2/2)^{\nu_0/2}}{\Gamma(\nu_0/2)}
 (\sigma^2)^{-(\nu_0/2+1)}e^{-\nu_0\sigma_0^2/(2\sigma^2)}.
\]
故联合先验的核为
\[
 \pi(\mu,\sigma^2)\propto(\sigma^2)^{-(\nu_0/2+3/2)}
 \exp\!\left[-\frac{\nu_0\sigma_0^2+\kappa_0(\mu-\mu_0)^2}{2\sigma^2}\right].
\]
\textbf{注意}：$\mu$ 的条件先验方差含 $\sigma^2$；不能把联合先验写成两个独立先验的乘积 $\pi(\mu)\pi(\sigma^2)$。

\block{12\quad 相乘、配方与参数更新}
\textbf{Step 3：似然乘先验，合并幂次和指数}
\[
 \pi(\mu,\sigma^2\mid x)\propto L(\mu,\sigma^2\mid x)\pi(\mu,\sigma^2).
\]
\[
 \begin{aligned}
 \pi(\mu,\sigma^2\mid x)
 &\propto(\sigma^2)^{-((\nu_0+n)/2+3/2)}\\
 &\quad\times\exp\!\bigg\{-\frac{1}{2\sigma^2}\Big[\nu_0\sigma_0^2+(n-1)s^2\\
 &\hspace{18mm}+\kappa_0(\mu-\mu_0)^2+n(\mu-\bar x)^2\Big]\bigg\}.
 \end{aligned}
\]
配方恒等式：
\[
 \begin{aligned}
 &\kappa_0(\mu-\mu_0)^2+n(\mu-\bar x)^2\\
 &=\kappa_n(\mu-\mu_n)^2+\frac{\kappa_0n}{\kappa_0+n}(\bar x-\mu_0)^2.
 \end{aligned}
\]
参数更新为
\[
 \boxed{\ \kappa_n=\kappa_0+n,\qquad\nu_n=\nu_0+n\ },
\]
\[
 \boxed{\ \mu_n=\frac{\kappa_0\mu_0+n\bar x}{\kappa_0+n}\ },
\]
\[
 \boxed{\ \nu_n\sigma_n^2=\nu_0\sigma_0^2+(n-1)s^2
 +\frac{\kappa_0n}{\kappa_0+n}(\bar x-\mu_0)^2\ }.
\]
\end{minipage}\hfill
\begin{minipage}[t]{.487\textwidth}
\block{13\quad 联合后验、条件后验与边际后验}
\textbf{Step 4：识别后验分布}
\[
 \pi(\mu,\sigma^2\mid x)\propto(\sigma^2)^{-(\nu_n/2+3/2)}
 \exp\!\left[-\frac{\nu_n\sigma_n^2+\kappa_n(\mu-\mu_n)^2}{2\sigma^2}\right].
\]
与联合先验同形，仍为正态—倒 Gamma 分布，并可分解为
\[
 \pi(\mu,\sigma^2\mid x)=\pi(\mu\mid\sigma^2,x)\pi(\sigma^2\mid x),
\]
\[
 \boxed{\ \mu\mid\sigma^2,x\sim N\!\left(\mu_n,\frac{\sigma^2}{\kappa_n}\right)\ },
\]
\[
 \boxed{\ \sigma^2\mid x\sim\IG\!\left(\frac{\nu_n}{2},\frac{\nu_n\sigma_n^2}{2}\right)\ }.
\]
$\mu\mid\sigma^2,x$ 是\textbf{条件后验}；$\sigma^2\mid x$ 是\textbf{边际后验}。
$\sigma_n^2$ 是更新后的超参数，不等同于 $\E(\sigma^2\mid x)$。

\textbf{$\mu$ 的边际后验}：把 $\sigma^2$ 积分掉。
\[
 \begin{aligned}
 \pi(\mu\mid x)
 &=\int_0^\infty\pi(\mu,\sigma^2\mid x)\,d\sigma^2\\
 &\propto\left[\nu_n\sigma_n^2+\kappa_n(\mu-\mu_n)^2\right]^{-(\nu_n+1)/2}\\
 &\propto\left[1+\frac1{\nu_n}
 \left(\frac{\mu-\mu_n}{\sigma_n/\sqrt{\kappa_n}}\right)^2\right]^{-(\nu_n+1)/2}.
 \end{aligned}
\]
因此
\[
 \boxed{\ \frac{\mu-\mu_n}{\sigma_n/\sqrt{\kappa_n}}\,\bigg|\,x\sim t_{\nu_n}\ }.
\]
这里的 $t_{\nu_n}$ 为标准 $t$ 分布；$\sigma_n^2/\kappa_n$ 为尺度的平方，\textbf{不是方差}。

\block{14\quad 本例的均值、方差与加权解释}
\[
 \E(\mu\mid x)=\mu_n\quad(\nu_n>1),
\]
\[
 \Var(\mu\mid x)=\frac{\nu_n}{\nu_n-2}\frac{\sigma_n^2}{\kappa_n}
 \quad(\nu_n>2),
\]
\[
 \E(\sigma^2\mid x)=\frac{\nu_n\sigma_n^2}{\nu_n-2}\quad(\nu_n>2).
\]
\[
 \mu_n=\frac{\kappa_0}{\kappa_0+n}\mu_0+\frac n{\kappa_0+n}\bar x.
\]
$\kappa_0$ 可解释为先验信息对应的样本量；$\kappa_n=\kappa_0+n$。
$\nu_n=\nu_0+n$；更新平方和由三部分相加：
\[
 \underbrace{\nu_0\sigma_0^2}_{\text{先验平方和}}
 +\underbrace{(n-1)s^2}_{\text{样本内平方和}}
 +\underbrace{\frac{\kappa_0n}{\kappa_0+n}(\bar x-\mu_0)^2}_{\text{两均值差异的附加项}}.
\]
\textbf{本例求方差可信上限}：由本例的边际后验，直接代入第 8 块：
\[
 \frac{\nu_n\sigma_n^2}{\sigma^2}\,\bigg|\,x\sim\chi^2(\nu_n),\qquad
 (\sigma^2)_U=\frac{\nu_n\sigma_n^2}{\chi^2_\gamma(\nu_n)}.
\]
\[
 P\bigl(\sigma^2\le(\sigma^2)_U\mid x\bigr)=1-\gamma.
\]
\end{minipage}
\end{document}
